\chapter*{能力清单}
\addcontentsline{toc}{chapter}{能力清单}

\begin{rubricbox}
\begin{itemize}
\tightlist
\item 能从晶体管开关、逻辑门、寄存器和时钟解释状态机为什么能成为 CPU 的基础。
\item 能手跑 X64S 教学子集的 \code{mov/add/jmp/syscall} 程序，说明 \code{RIP}、\code{RSP}、寄存器、\code{RFLAGS}、\code{CPL} 如何变化。
\item 能从轮询设备推导出中断，从死循环推导出定时器抢占，从同权代码推导出特权级和系统调用。
\item 能画出一条 x86-64 load/store 指令经过 RIP、寄存器、ALU、MMU、TLB、cache 和异常入口的路径。
\item 能区分 ISA、微架构和 ABI，并说明 OS 哪些代码依赖 ISA，哪些性能现象来自微架构。
\item 能解释总线、MMIO、DMA、IOMMU、中断控制器和设备描述符环怎样共同完成一次 I/O。
\item 能说明特权级限制了哪些指令和状态，为什么系统调用必须是受控陷入而不是普通函数调用。
\item 能把一次系统调用拆成用户态入口、内核对象、权限检查、状态更新和返回值。
\item 能解释线程为何没有推进：等 CPU、等锁、等页面、等 I/O、等信号还是等子进程。
\item 能解释从复位向量到 \code{kernel\_main} 的启动阶段和每个阶段禁止事项。
\item 能说明中断上半部、下半部、定时器 tick、抢占点和上下文保存。
\item 能判断自旋锁、mutex、semaphore、RCU 的上下文限制，并用锁顺序图解释死锁。
\item 能解释 \code{exec} 的 ELF 验证、段映射、用户栈构造、动态解释器和提交点。
\item 能手算一次虚拟地址翻译，并区分合法缺页、权限错误和非法地址。
\item 能把一次 x86-64 \#PF 拆成 CR2、error code、VMA 查找、PTE 判定、COW/page cache 修复、TLB 失效和返回重试。
\item 能画出 VMA、PTE、page cache、inode 的关系，并区分 \code{MAP\_SHARED} 与 \code{MAP\_PRIVATE}。
\item 能比较 buddy、slab、kmalloc size class 的用途、复杂度和碎片成本。
\item 能画出驱动描述符环、DMA 映射、completion 和 buffer 生命周期。
\item 能说明 bio、request、request queue、flush、FUA 和 barrier 如何共同影响 I/O 顺序。
\item 能画出 fd、open file description、inode、page cache 和设备请求之间的对象关系。
\item 能解释 inode、目录项、空闲块 bitmap、VFS 和 page cache 如何协作。
\item 能解释 socket 作为 fd 的生命周期，说明 bind、listen、accept、connect、send、recv 的等待语义。
\item 能分析 capability、namespace、cgroup、seccomp 如何共同形成最小权限边界。
\item 能区分 errno、信号、WARN、oops、panic，并为系统调用写出失败回滚路径。
\item 能为可靠写出设计 prepared、synced、renamed、manifest committed 的提交状态机。
\item 能写出包含现象、假设、证据、反证和未验证边界的 OS 调试报告。
\end{itemize}
\end{rubricbox}
